\section{偶関数と奇関数}\label{節:奇偶の対称性}

\begin{definition}{偶関数と奇関数}
原点に関して対称な範囲で定義された関数 $f$ について,すべての $x$ で
\[
f(-x)=f(x)
\]
が成り立つとき,\,$f$ を偶関数という.また,すべての $x$ で
\[
f(-x)=-f(x)
\]
が成り立つとき,\,$f$ を奇関数という.
\end{definition}

偶関数のグラフは $y$ 軸について対称であり,奇関数のグラフは原点について対称である.\,$x^2$ と $\cos x$ は偶関数,\,$x^3$ と $\sin x$ は奇関数である.ただし,式の中に $\sin$ があれば奇関数になるわけではない.\,$\sin(x^2)$ では $x$ を $-x$ に変えても $x^2$ が変わらないので偶関数である.式全体に $-x$ を代入することが,確実な判定になる.

区間が原点について対称なとき,この対称性を積分に使えないだろうか.$a>0$ とし,\,$[-a,a]$ での積分を,区間 $[0,a]$ での積分に帰着させることを目標とする.

$[-a,a]$ の積分を $0$ で分ける.左半分で $x=-t$ と置くと,\,$x=-a,0$ に対して $t=a,0$ であり,\,$\dd x=-\dd t$ だから,
\[
\begin{aligned}
\int_{-a}^{0}f(x)\dd x
 &=\int_a^0 f(-t)(-\dd t)\\
 &=\int_0^a f(-t)\dd t
\end{aligned}
\]
となる.積分変数の文字を $x$ に戻して右半分と加えれば,
\begin{equation}\label{式:左右の積分}
\int_{-a}^a f(x)\dd x
 =\int_0^a\{f(x)+f(-x)\}\dd x
\end{equation}
を得る.これは偶関数でも奇関数でもない場合にも使える式である.偶関数なら中括弧が $2f(x)$ になり,奇関数なら $0$ になる.

\begin{theorem}{対称区間での偶関数・奇関数の積分}
$a>0$ とする.\,$f$ が偶関数なら
\[
\int_{-a}^a f(x)\dd x=2\int_0^a f(x)\dd x,
\]
$f$ が奇関数なら
\[
\int_{-a}^a f(x)\dd x=0
\]
である.
\end{theorem}

図\ref{図:奇偶の積分}の右図では,左右の部分の幾何学的な面積は等しい.しかし左側の定積分は負,右側は正である.奇関数の積分が $0$ という結論は,図形の面積が $0$ という意味ではない.面積を求めるなら絶対値を積分する必要がある.

\begin{figure}[htbp]
\centering
\begin{tikzpicture}[scale=1.05,>=Stealth]
\begin{scope}[xshift=-3.1cm]
\fill[blue!12] plot[domain=-1.35:1.35,samples=50](\x,{0.65*\x*\x})--(1.35,0)--(-1.35,0)--cycle;
\draw[->](-1.7,0)--(1.8,0)node[right]{$x$};
\draw[->](0,-0.2)--(0,1.7)node[above]{$y$};
\draw[blue!70!black,thick] plot[domain=-1.45:1.45,samples=60](\x,{0.65*\x*\x});
\draw[dashed](-1.35,0)--(-1.35,1.185);
\draw[dashed](1.35,0)--(1.35,1.185);
\node[below]at(-1.35,0){$-a$};\node[below]at(1.35,0){$a$};
\node at(0,-0.65){偶関数：左右を加える};
\end{scope}
\begin{scope}[xshift=1.8cm]
\fill[blue!12](0,0)--plot[domain=0:1.35,samples=40](\x,{0.55*\x*\x*\x})--(1.35,0)--cycle;
\fill[orange!18](0,0)--plot[domain=0:-1.35,samples=40](\x,{0.55*\x*\x*\x})--(-1.35,0)--cycle;
\draw[->](-1.7,0)--(1.8,0)node[right]{$x$};
\draw[->](0,-1.55)--(0,1.7)node[above]{$y$};
\draw[blue!70!black,thick] plot[domain=-1.4:1.4,samples=60](\x,{0.55*\x*\x*\x});
\draw[dashed](-1.35,0)--(-1.35,-1.353);
\draw[dashed](1.35,0)--(1.35,1.353);
\node[above]at(-1.35,0){$-a$};\node[below]at(1.35,0){$a$};
\node at(0,-1.95){奇関数：符号付きの量が打ち消す};
\end{scope}
\end{tikzpicture}
\caption{対称性と,定積分の符号の関係.曲線はそれぞれ $x^2$ 型,\,$x^3$ 型である.}\label{図:奇偶の積分}
\end{figure}

\begin{bookpoint}
\textbf{POINT}\par
$x$ を $-x$ に変えて各項の偶奇を調べ,奇関数の項は $0$ とし,偶関数の項は $[0,a]$ での積分の2倍にする.
\end{bookpoint}

\begin{example}{偶奇による分解}
次の定積分を求めよ.
\[
I=\int_{-2}^{2}\left(x^3\e^{-x^2}+2x^2+1\right)\dd x.
\]
\end{example}

\paragraphbf{方針}
被積分関数を $x^3\e^{-x^2}$ と $2x^2+1$ に分け,それぞれの偶奇を調べる.

\paragraphbf{解答}
$I=\frac{44}{3}$

\paragraphbf{解法}
$x$ を $-x$ に変えると,\,$x^3$ の符号だけが変わり,\,$\e^{-x^2}$ は変わらない.したがって積 $x^3\e^{-x^2}$ は奇関数で,\,$[-2,2]$ での積分は $0$ である.残った $2x^2+1$ は偶関数だから,
\[
\begin{aligned}
I&=2\int_0^2(2x^2+1)\dd x\\
 &=2\left[\frac{2x^3}{3}+x\right]_0^2\\
 &=2\left(\frac{16}{3}+2\right)\\
 &=\frac{44}{3}
\end{aligned}
\]
である.\,$x^3\e^{-x^2}$ の原始関数は求めていない.また,式全体が偶関数か奇関数かを決める必要もなかった.和を分けて,対称性が分かる部分にだけ使えばよい.

\paragraphbf{偶関数部分と奇関数部分}
さらに,一般の $f$ に対して
\[
f_{\mathrm e}(x)=\frac{f(x)+f(-x)}{2},\qquad
f_{\mathrm o}(x)=\frac{f(x)-f(-x)}{2}
\]
と置くと,\,$f=f_{\mathrm e}+f_{\mathrm o}$ であり,前者は偶関数,後者は奇関数になる.よって対称区間の定積分に寄与するのは $f_{\mathrm e}$ の部分だけである.実際,式\eqref{式:左右の積分}の中括弧は $f(x)+f(-x)=2f_{\mathrm e}(x)$ であり,この式は $f$ から偶関数の部分を取り出して $2\int_0^a f_{\mathrm e}(x)\dd x$ を計算していたのである.
